% Generated by ChimCanhCut Research Platform
% Template: IEEEtran — IEEE Transactions on Neural Networks and Learning Systems / TPAMI / TKDE
% Topic: Federated Learning under Concept Drift: A Unified Perspective on Drift-Aware Learning in Distributed Non-Stationary Environments

\documentclass[journal]{IEEEtran}

\usepackage[utf8]{inputenc}
\usepackage{cite}
\usepackage{amsmath,amssymb}
\usepackage{graphicx}
\usepackage{booktabs}
\usepackage{hyperref}
\usepackage{url}
\usepackage{multirow}
\usepackage{array}

\title{Federated Learning under Concept Drift: A Unified Perspective on Drift-Aware Learning in Distributed Non-Stationary Environments}
\author{[Author Names]}

\begin{document}
\maketitle

\begin{abstract}
% TODO: Enter your abstract via the Paper Metadata panel in the UI.
% Required: 150-250 words covering motivation, methodology, key findings, contributions.
\end{abstract}

\begin{IEEEkeywords}Federated Learning, Concept Drift, Non-IID Data, Distribution Shift\end{IEEEkeywords}

\section{Introduction}
\label{sec:introduction}

Federated Learning (FL) has emerged as a prominent paradigm for privacy-preserving collaborative machine learning \cite{Collins2025}. It enables distributed clients, such as Internet of Things (IoT) devices and edge nodes, to jointly train global models without sharing raw local data [CITE:34, CITE:104]. However, real-world distributed environments are inherently dynamic. While existing FL literature extensively addresses static non-IID (Independent and Identically Distributed) data challenges [CITE:90, CITE:91, CITE:92], an increasingly prevalent issue is concept drift—the phenomenon where the underlying data distribution changes over time \cite{Martin2025}. 

In distributed networks, concept drift manifests not only temporally but also spatially across heterogeneous clients, leading to distributed concept drift [CITE:14, CITE:22]. When traditional FL aggregation methods, such as FedAvg, are deployed in these non-stationary environments, they often experience significant model degradation and catastrophic forgetting [CITE:48, CITE:78]. This challenge is further compounded in IoT and edge deployments, where strict constraints on communication bandwidth, computational capacity, and battery life restrict the feasibility of continuous retraining or complex model maintenance [CITE:73, CITE:101, CITE:122].

To address these challenges, a rapidly expanding body of literature has proposed drift-aware federated learning frameworks. To provide a structured understanding of this field, this survey introduces a comprehensive taxonomy that categorizes existing solutions across three primary dimensions: adaptation strategies, network topologies, and the locus of adaptation. 

First, regarding *adaptation strategies*, approaches generally divide into active and passive methods. Active approaches rely on explicit drift detection mechanisms—such as statistical triggers, loss-based monitoring, or dual-schedulers—to initiate reactive model updates [CITE:4, CITE:27, CITE:34, CITE:45]. Conversely, passive strategies employ continuous adaptation techniques, such as dynamic learning rate adjustments or regularization, to maintain model plasticity during drift phases without requiring computationally demanding statistical detectors [CITE:17, CITE:31, CITE:40, CITE:44]. 

Second, *network topologies* dictate how drift is managed across the communication infrastructure. Solutions vary from traditional centralized FL, where a single server orchestrates drift adaptation, to decentralized (peer-to-peer) and hierarchical architectures that localize drift responses to reduce communication latency and mitigate single points of failure [CITE:55, CITE:88]. 

Third, the *locus of adaptation* identifies where the primary drift mitigation occurs. Client-side adaptation focuses on localized learning and personalized models [CITE:11, CITE:14]; server-side adaptation involves intelligent aggregation and global model clustering [CITE:22, CITE:48]; and hybrid approaches distribute the adaptive workload across both clients and the server to balance personalization with global generalization [CITE:60, CITE:105].

Within this taxonomic framework, the structural maintenance of models under drift varies significantly. While some centralized frameworks adapt a single global model \cite{Casado2021}, others argue that distributed concept drift necessitates multi-model or clustered approaches to capture divergent client distributions effectively [CITE:11, CITE:14, CITE:22]. To manage the inherent trade-off between accommodating new concepts and retaining past knowledge, there is a growing trend toward integrating Federated Continual Learning (FCL) techniques. Methods such as experience replay and elastic weight consolidation treat concept drift as a continuous learning problem [CITE:112, CITE:117, CITE:128]. However, deploying FCL in edge environments requires careful optimization, as the memory and processing demands of retaining historical knowledge must be directly balanced against the hardware limitations of IoT devices [CITE:73, CITE:122].

Despite these advancements, several critical research gaps remain. First, there is a notable absence of unified, standardized benchmarks for evaluating non-stationary federated learning. The literature frequently relies on fragmented datasets or simplistic synthetic adaptations, limiting rigorous cross-paper comparative analyses. Second, active drift detection mechanisms often make strong assumptions about the computational capacity of resource-constrained edge devices to continuously evaluate models on new samples or maintain complex statistical triggers. Third, multi-model and clustered approaches designed to handle distributed concept drift face scalability challenges, including increased storage and communication overheads as the number of models grows over time. Additionally, there is a theoretical gap regarding how central servers, which lack access to local data distributions, can accurately identify specific drifting clients during aggregation. Finally, the intersection of model fairness and group-specific distributed concept drift remains largely unexplored, presenting an opportunity to develop frameworks that ensure equitable outcomes across sensitive attributes as data distributions evolve [CITE:15, CITE:23].

This survey provides a unified perspective on drift-aware learning in distributed non-stationary environments \cite{Mahdi2025}. By systematically reviewing the foundations, innovations, and existing literature of federated learning under concept drift, we aim to consolidate current knowledge through our proposed three-dimensional taxonomy. Ultimately, this work critically analyzes the trade-offs within adaptation strategies, network topologies, and loci of adaptation, outlining concrete future research directions to build robust, scalable, and fair federated systems.

\begin{figure}[!t]
\centering
\includegraphics[width=\columnwidth]{figures/year_distribution}
\caption{Publication year distribution of included papers.}
\label{fig:year_dist}
\end{figure}

\section{Background and Related Work}
\label{sec:background}

Federated Learning (FL) has emerged as a standard paradigm for privacy-preserving collaborative intelligence, enabling distributed clients to train machine learning models without sharing raw data \cite{Collins2025}. While FL effectively addresses data privacy, it encounters well-documented challenges with spatial data heterogeneity, commonly referred to as the Non-IID (Independent and Identically Distributed) data problem [CITE:90, 91, 92]. Extensive research has explored techniques such as clustered federated learning, personalized federated learning, and data augmentation to mitigate the effects of skewed data distributions across clients [CITE:94, 99, 102, 105, 107]. However, traditional FL formulations generally assume a static environment where local data distributions, although heterogeneous, remain stationary over time.

In real-world applications, particularly those involving streaming data, environments are rarely static. The statistical properties of the target variables or input data often change over time, a phenomenon known as concept drift [CITE:75, 77, 79]. In centralized machine learning, concept drift is traditionally managed through continuous monitoring and adaptation, utilizing distribution-based, performance-based, or unsupervised drift detectors [CITE:71, 78, 81]. Systematic reviews of streaming data classification highlight various active detection algorithms and passive adaptation strategies, such as dynamic ensemble selection and incremental learning [CITE:55, 63, 74, 83].

When concept drift occurs within a federated network, the system faces compounding challenges of both spatial heterogeneity and temporal non-stationarity, leading to "distributed concept drift" [CITE:14, 22]. Empirical studies indicate that concept drift in FL can significantly degrade global model performance if not explicitly addressed [CITE:33, 42]. Furthermore, managing drift in a distributed context is complicated by an inherent information asymmetry: the central server lacks direct access to local data distributions, limiting its ability to identify which specific clients are experiencing drift during the aggregation phase.

To navigate these challenges, the literature has proposed various methodologies, which this survey categorizes into a unified taxonomy based on three core dimensions: *adaptation strategies*, *network topologies*, and the *locus of adaptation*. First, *adaptation strategies* differentiate between active approaches, which rely on explicit drift detection mechanisms (e.g., statistical triggers) to initiate reactive model updates [CITE:34, 43, 45], and passive strategies, which utilize continuous dynamic adjustments (e.g., learning rate modulation) to maintain model plasticity without requiring explicit detectors [CITE:17, 31, 40, 44]. Second, *network topologies* define the structural response to drift, ranging from maintaining a single evolving global model to decentralized or clustered multi-model architectures designed to capture divergent client distributions as new concepts emerge [CITE:1, 11, 14, 21, 22]. Finally, the *locus of adaptation* specifies whether the primary drift mitigation occurs on the server-side (e.g., via specialized aggregation weights and momentum) or the client-side (e.g., via local model regularization or personalized fine-tuning) [CITE:27, 38, 52].

A critical aspect of client-side adaptation involves mitigating catastrophic forgetting—the loss of previously acquired knowledge when learning new concepts from non-stationary data. Consequently, Continual Learning (CL) techniques are increasingly integrated into FL frameworks, treating concept drift as a continuous learning problem [CITE:48, 112, 117]. Federated Continual Learning (FCL) frameworks employ strategies like prioritized experience replay, elastic weight consolidation, and knowledge distillation to retain prior concepts while adapting to new ones [CITE:118, 122, 128]. However, integrating these sophisticated FCL mechanisms introduces substantial computational and memory demands. This creates a logical tension when deploying drift-aware systems in resource-constrained IoT and edge environments [CITE:88, 101, 104]. Active drift detection and continuous replay mechanisms often assume that edge devices can continuously evaluate complex statistical triggers or store historical samples, which is frequently impractical \cite{ren2025}. As a result, recent research emphasizes co-optimizing drift adaptation with communication efficiency and hardware latency constraints [CITE:34, 45, 50, 60].

Despite these advancements, the field of drift-aware federated learning remains fragmented \cite{Mahdi2025}. A notable limitation is the absence of standardized benchmarks, which forces researchers to rely on proprietary datasets or simplistic synthetic adaptations, complicating rigorous cross-paper comparative analyses of non-stationary FL algorithms. Furthermore, emerging challenges, such as ensuring fairness under group-specific distributed concept drift, remain largely unexplored [CITE:15, 23]. By explicitly structuring the literature across adaptation strategies, network topologies, and the locus of adaptation, this survey provides a unified perspective. This framework aims to synthesize existing methodologies, clarify the operational trade-offs between different architectural choices, and guide the development of scalable, resource-efficient FL systems in dynamic environments.

\section{Problem Formulation}
\label{sec:problem_formulation}

The problem of Federated Learning (FL) under concept drift emerges at the intersection of spatial data heterogeneity and temporal non-stationarity. In a standard FL formulation, a central server or a network of nodes coordinates $N$ distributed clients to minimize a global empirical risk function without directly accessing local datasets. Traditional FL implicitly assumes that the underlying data distributions remain stationary over time. However, in real-world distributed environments, such as Internet of Things (IoT) and edge networks, this assumption is frequently violated [CITE:34, 39, 73]. 

Concept drift introduces temporal heterogeneity, meaning the joint probability distribution of the input features and target labels for a given client $i$ at time $t$, denoted as $P_i^t(X, Y)$, changes over time such that $P_i^t(X, Y) \neq P_i^{t+1}(X, Y)$ [CITE:33, 42, 79]. When this temporal evolution occurs across multiple distributed nodes with non-independent and identically distributed (non-IID) data [CITE:90, 91, 92], it manifests as *distributed concept drift* [CITE:14, 22]. Formulating a comprehensive solution to distributed concept drift requires a structured approach to several interconnected challenges. To provide a unified perspective, we formulate the problem space along three primary dimensions: (1) *drift adaptation strategies*, (2) *the primary locus of adaptation*, and (3) *network topologies*.

The first dimension, **drift adaptation strategies**, addresses how the system identifies and responds to non-stationarity. A primary dichotomy exists between active and passive adaptation. Active detection paradigms explicitly monitor statistical triggers, loss metrics, or data distributions to identify drift and initiate reactive model updates [CITE:34, 43, 75, 78]. While frameworks utilizing explicit detection (e.g., FedConD, FLARE) can precisely trigger retraining or model transfers [CITE:34, 39, 45], they often rely on the computationally demanding assumption that resource-constrained edge devices can continuously evaluate statistical triggers without incurring prohibitive latency overheads [CITE:73, 88]. Conversely, passive adaptation strategies bypass explicit detection, utilizing mechanisms such as dynamic learning rates or regularization to continuously maintain model plasticity [CITE:17, 31, 40, 44]. Furthermore, as local streams evolve, models are susceptible to catastrophic forgetting. Consequently, adaptation strategies frequently intersect with Federated Continual Learning (FCL), treating drift adaptation as a continuous process that must balance plasticity (learning new concepts) and stability (retaining prior knowledge) under strict resource constraints [CITE:48, 112, 117, 122].

The second dimension concerns the **locus of adaptation**, delineating whether drift mitigation is primarily managed at the client or server level. Client-side adaptation focuses on localized mechanisms, such as personalized learning rates or local memory buffers, to handle drift before aggregation [CITE:17, 112]. Server-side adaptation involves adjusting aggregation weights, managing model clusters, or routing specialized models to specific clients [CITE:11, 20]. This dimension directly impacts architectural formulations. Maintaining a single global model is often insufficient when distributed drift causes clients to diverge into distinct conceptual spaces. To address this, multi-model and clustered FL approaches (e.g., FedDrift, FedCCFA) dynamically group clients experiencing similar drifts [CITE:14, 22]. However, these formulations introduce scalability challenges, including the management of multiple models, complex merging heuristics, and increased communication overhead. Additionally, privacy constraints render the server largely "blind" to local data distributions, complicating the accurate identification of drifting clients during server-side aggregation [CITE:4, 27].

The third dimension involves **network topologies**, which dictate the flow of knowledge and the coordination of drift adaptation. While the majority of drift-aware FL formulations assume a centralized, client-server topology, this structure can create communication bottlenecks and delayed responses during widespread, rapid concept drift. Decentralized (peer-to-peer) topologies offer an alternative formulation where clients share models directly with neighbors, potentially allowing for faster, localized adaptation to regional drifts without relying on a central coordinator [CITE:55, 62]. Hierarchical topologies further bridge this gap by utilizing edge servers to manage localized drift clusters before global aggregation [CITE:68, 71]. The choice of topology fundamentally alters how distributed concept drift is detected and mitigated, influencing both communication efficiency and adaptation speed.

Finally, emerging formulations must account for the broader socio-technical implications of non-stationarity, such as group-specific distributed concept drift. When drift disproportionately affects certain demographic groups or sensitive attributes, it raises important fairness and robustness considerations that must be integrated into the learning objective [CITE:15, 23]. 

In summary, formulating the problem of drift-aware federated learning requires navigating these three dimensions—adaptation strategies, locus of adaptation, and network topologies—to optimize communication efficiency, mitigate catastrophic forgetting, and dynamically adapt to spatio-temporal heterogeneity within the computational limits of distributed networks.

\section{Methodology}
\label{sec:methodology}

\#\# Methodology

To ensure a rigorous, transparent, and reproducible synthesis of the literature concerning federated learning (FL) in non-stationary environments, this systematic survey was conducted following the Preferred Reporting Items for Systematic Reviews and Meta-Analyses (PRISMA) 2020 guidelines. The methodology encompasses a well-defined search strategy, strict eligibility criteria, a multi-stage screening process, and structured data extraction protocols.

\#\#\# 1. Search Strategy
A comprehensive literature search was executed to identify relevant studies at the intersection of federated learning and concept drift. The search targeted major academic databases and preprint repositories, specifically **Semantic Scholar** and **arXiv**, with supplementary indexing planned for **IEEE Xplore (pending)** to ensure exhaustive coverage of both computer science literature and broader engineering applications. 

The search period was strictly defined as **2015 to 2026**. The year 2015 was selected as the foundational baseline to capture early decentralized machine learning paradigms immediately preceding the formal introduction of federated learning, while extending to 2026 ensures the inclusion of the most recent, cutting-edge advancements and accepted forthcoming publications. 

To capture the diverse terminology used across different subfields, a robust set of **45 distinct search queries** was constructed. These queries were formulated using Boolean operators to intersect core concepts from three primary domains: 
*   *Distributed Learning:* ("federated learning" OR "collaborative learning" OR "decentralized learning")
*   *Non-Stationary Environments:* ("concept drift" OR "data shift" OR "non-stationary" OR "temporal shift")
*   *Adaptive Mechanisms:* ("drift-aware" OR "continuous learning" OR "dynamic adaptation")

\#\#\# 2. Inclusion and Exclusion Criteria
To maintain the scientific rigor and topical focus of this survey, retrieved records were subjected to predefined inclusion and exclusion criteria. 

**Inclusion Criteria:**
1.  **Published in peer-reviewed venues or arXiv preprints:** Justified by the exceptionally rapid pace of innovation in artificial intelligence, where preprints often contain highly relevant, state-of-the-art methodologies before formal peer-review cycles conclude.
2.  **Directly or adjacently related to the survey topic:** Studies must address either explicit concept drift in FL or related non-stationary challenges (e.g., continual learning, domain adaptation) that inform drift-aware FL design.
3.  **Published after 2015:** Justified by the temporal scope of modern federated learning architectures.
4.  **Abstract available:** Essential for the initial semantic screening phases.

**Exclusion Criteria:**
1.  **Duplicate entries:** To prevent redundant analysis and skewed quantitative assessments.
2.  **Non-English papers:** Excluded to ensure accurate interpretation and synthesis of complex algorithmic concepts.
3.  **Papers without an accessible abstract:** Excluded as they prevent efficient preliminary relevance evaluation.
4.  **Papers published before 2015:** Excluded to filter out legacy distributed computing models that do not align with modern privacy-preserving FL constraints.

\#\#\# 3. Study Selection and PRISMA Flow
The study selection process followed a structured, multi-stage PRISMA pipeline:

*   **Identification:** The execution of the 45 search queries across the selected databases yielded a total of **129 records**. 
*   **Screening:** All **129 records** underwent a two-tiered screening process based on titles and abstracts. First, a rule-based filtering mechanism was applied to enforce baseline criteria (e.g., date ranges, language, and duplication), resulting in the exclusion of **19 records**. Subsequently, an advanced Large Language Model (LLM)-assisted screening protocol was utilized to evaluate the semantic alignment of the remaining abstracts with the core topic. This LLM-assisted phase identified and excluded **3 records** that used matching keywords but lacked contextual relevance to federated non-stationary environments.
*   **Eligibility:** The remaining **107 records** advanced to the eligibility phase, where full texts were retrieved and thoroughly assessed. Because the preceding multi-tiered screening was highly precise, **0 full-text records were excluded** at this stage.
*   **Included:** Ultimately, **107 studies** met all criteria and were included in the final qualitative and quantitative synthesis. 

\#\#\# 4. Quality Assessment
Rather than relying exclusively on journal impact factors, the quality and relevance of the included papers were assessed based on methodological soundness, experimental validation, and topological relevance to the survey's core theme. To structure the narrative of the review, the **107 included studies** were stratified into three distinct relevance categories:
*   **Direct (53 studies):** Papers that explicitly propose, analyze, or benchmark drift-aware mechanisms within federated learning architectures.
*   **Adjacent (53 studies):** Papers that address highly related challenges—such as statistical heterogeneity (Non-IID data), federated continual learning, or distributed domain adaptation—which provide adaptable techniques for managing concept drift.
*   **Foundational (1 study):** A seminal paper utilized to establish the core theoretical definitions and mathematical grounding of concept drift in distributed environments.

\#\#\# 5. Data Extraction
To facilitate a unified perspective on the literature, a standardized data extraction protocol was applied to the 107 included studies. For each paper, the following key variables were systematically extracted:
1.  **Drift Typology:** The specific nature of the concept drift addressed (e.g., sudden, gradual, incremental, recurring, or feature/label shift).
2.  **FL Architecture:** The underlying distributed topology (e.g., cross-device, cross-silo, decentralized/peer-to-peer, or hierarchical).
3.  **Algorithmic Mechanisms:** The core strategies utilized for drift detection (e.g., statistical testing, loss-monitoring) and drift adaptation (e.g., dynamic weighting, model personalization, knowledge distillation).
4.  **Privacy and Communication Constraints:** How the proposed drift-aware methods balance adaptation speed with communication overhead and data privacy guarantees.
5.  **Evaluation Frameworks:** The specific datasets (synthetic and real-world), simulation environments, and performance metrics used to validate the proposed methodologies. 

This structured extraction process ensures a comprehensive synthesis of current capabilities, enabling the identification of critical research gaps in drift-aware federated learning.

\begin{figure}[!t]
\centering
\includegraphics[width=0.85\columnwidth]{figures/prisma_flow}
\caption{PRISMA 2020 flow diagram of the systematic literature search.}
\label{fig:prisma}
\end{figure}

\section{Taxonomy}
\label{sec:taxonomy}

To systematically categorize the diverse methodologies addressing non-stationary data in distributed environments, we propose a multi-dimensional taxonomy. This taxonomy provides a unified perspective by classifying existing literature along three intersecting axes: Drift Adaptation Strategy, Primary Locus of Adaptation, and Network Topology. By examining how algorithms detect or absorb drift (Strategy), where the computational burden of this adaptation resides (Locus), and how the communication infrastructure facilitates model updates (Topology), we can comprehensively evaluate the design trade-offs of current drift-aware federated systems.

**1. Drift Adaptation Strategy**
The approach to managing concept drift represents a primary point of categorization in the literature, broadly divided into Active Detection \& Reactive, Passive Adaptation, and Clustering \& Multi-Model strategies.

*Active Detection \& Reactive:* Frameworks in this category utilize explicit mechanisms to identify concept drift and subsequently trigger model updates. These approaches typically rely on statistical tests, dual-schedulers, or historical performance metrics to initiate retraining. For instance, the FLARE framework employs conditional transfer of training data and model deployment based on training behavior and inference statistics \cite{Chow2023}, \cite{Chow2023a}. Similarly, FLAME utilizes dynamic thresholding for drift detection \cite{Mavromatis2024}, while FL-MalDrift integrates on-device statistical drift detection algorithms such as ADWIN and DDM \cite{Patel2025}. While effective at isolating drift points, active detection often requires clients to continuously evaluate models on new samples or maintain statistical tracking. This can present a computational and latency challenge for resource-constrained edge and Internet of Things (IoT) devices.

*Passive Adaptation:* In contrast to explicit detection, passive strategies favor continuous, lightweight adjustments to maintain model plasticity during drift phases, thereby mitigating the computational overhead associated with active triggers. A prevalent approach involves dynamic or adaptive learning rates based on individual training loss, as demonstrated in Adaptive-FedAVG \cite{Canonaco2021a}, \cite{Canonaco2021} and various client-side adaptation algorithms \cite{Saile2024a}, \cite{Saile2024}. Additionally, Continual Learning (CL) techniques are frequently integrated into passive FL frameworks to address catastrophic forgetting. These methods treat concept drift as an ongoing learning process rather than a discrete anomaly detection task, utilizing regularization or replay buffers (e.g., CDA-FedAvg \cite{Casado2021a}, \cite{Casado2021} and Federated Continual Learning frameworks \cite{Hamedi2025}, \cite{Benameur2025}).

*Clustering \& Multi-Model:* This approach addresses the variation in data distributions across the network. While standard FL maintains a single evolving global model, multi-model frameworks posit that distributed concept drift often necessitates maintaining distinct models to capture divergent client distributions. Algorithms such as FedDrift \cite{Jothimurugesan2022}, \cite{Jothimurugesan2022a} cluster clients based on empirical risk and assign specialized models to handle different emerging concepts. Other examples include FedCCFA, which utilizes classifier clustering and feature alignment \cite{Chen2024}, \cite{Chen2024a}, and FedPLC, which implements dynamic cluster adaptation to group clients experiencing similar drift patterns \cite{Zhou2026}, \cite{Zhou2026a}.

**2. Primary Locus of Adaptation**
The responsibility for managing concept drift can be localized to different components of the federated architecture. The choice of locus directly impacts communication overhead, privacy guarantees, and the system's responsiveness to local versus global drift.

*   *Client-Side Adaptation:* In this paradigm, the mechanisms for drift handling are executed locally on edge devices. This approach ensures immediate responsiveness to local data shifts and strict privacy preservation, as raw statistical triggers or replay buffers never leave the device. Examples include client-side passive adaptation using dynamic learning rates \cite{Saile2024a}, \cite{Saile2024} and active local model loss-based detection \cite{Rahman2025}, \cite{Rahman2025a}. However, client-side approaches are inherently limited by the local device's computational constraints and lack a global perspective, potentially leading to overfitting on transient local noise.
*   *Server-Side Adaptation:* Here, the central server orchestrates drift management, leveraging its broader view of the network to distinguish between localized anomalies and global concept drift. Server-side strategies often involve dynamic aggregation weights, where the server discounts updates from clients experiencing severe drift or prioritizes those with novel concepts. This locus also handles efficient storage management for exemplar-based online continual learning, reducing the memory burden on clients \cite{Zhou2024}, \cite{paper62}. While this centralizes the computational load, it can introduce latency in responding to rapid, localized shifts.
*   *Hybrid / Joint Adaptation:* To balance the trade-offs of strictly client- or server-side approaches, hybrid frameworks distribute the adaptation workload. Typically, clients perform lightweight local monitoring and send drift signals or specific metadata to the server, which then adjusts the global aggregation strategy accordingly. FLARE, for example, uses a lightweight dual-scheduler that requires cooperation between local training behavior and global deployment \cite{Chow2023}, \cite{Hinder2024}. This coordinated effort allows for nuanced drift management that is both globally informed and locally responsive.

**3. Network Topology**
The underlying communication architecture significantly influences how drift is managed, particularly concerning resource efficiency, scalability, and the speed of adaptation in non-stationary environments.

*   *Standard Client-Server FL:* Traditional federated setups, where all clients communicate directly with a single central server, provide a straightforward mechanism for global model updates. However, in the presence of rapid concept drift, this star topology can suffer from communication bottlenecks, as frequent global aggregations are required to keep the model current \cite{Abozaid2025}. Consequently, adaptation strategies in this topology often focus on communication-efficient drift signaling \cite{Gutierrez2024}.
*   *Hierarchical / Edge-Assisted FL:* Designed to reduce communication overhead and latency, these topologies introduce intermediate edge servers. This structure is particularly advantageous for managing spatial concept drift, where different geographic regions experience distinct data shifts. Edge servers can perform localized drift adaptation and aggregate regional concepts before communicating with the central cloud. Examples include Hierarchical Bidirectional Averaging (HBA) for tree-topology networks \cite{Hemadasa2025} and PFL-DRIFT, which utilizes a two-level hierarchical aggregation paradigm to isolate and adapt to regional drift more rapidly than a standard centralized server \cite{Aziz2025a}, \cite{Aziz2025}.
*   *Decentralized and Asynchronous FL:* To mitigate straggler effects and handle heterogeneous latency—which are exacerbated when clients require varying computational time to adapt to new concepts—asynchronous frameworks allow clients to update the global model independently. Frameworks like FedConD utilize asynchronous updates to ensure that clients experiencing drift can adapt without waiting for a global synchronization step \cite{Chen2021}, \cite{Chen2021a}. Furthermore, fully decentralized (peer-to-peer) cooperative frameworks have emerged to handle continual learning without a central aggregator. By utilizing gossip protocols and localized neighborhood aggregation, these topologies allow customized models to evolve organically across the network, providing robust adaptation to non-stationary data even in the absence of centralized coordination \cite{Yang2025}, \cite{Lanza2025}, \cite{paper130}.

\begin{table}[!t]
\renewcommand{\arraystretch}{1.3}
\caption{Proposed Multi-Dimensional Taxonomy}
\label{tab:taxonomy}
\centering
\begin{tabular}{p{3cm}p{9cm}}
\toprule
\textbf{Dimension} & \textbf{Categories} \\
\midrule
\textbf{Drift Adaptation Strategy} & Passive Adaptation, Active Detection \& Reactive, Clustering \& Multi-Model \\
\textbf{Network Topology} & Standard Client-Server FL, Hierarchical / Edge-Assisted FL, Decentralized / Asynchronous FL, Non-Federated Data Streams \\
\textbf{Primary Locus of Adaptation} & Client-Side, Server-Side, Hybrid / Joint, Not Applicable \\
\bottomrule
\end{tabular}
\end{table}

\begin{figure}[!t]
\centering
\includegraphics[width=\columnwidth]{figures/taxonomy_heatmap}
\caption{Heatmap of paper distribution across taxonomy dimensions.}
\label{fig:taxonomy_heatmap}
\end{figure}

\section{Literature Review}
\label{sec:literature_review}

The study of concept drift in Federated Learning (FL) has evolved from addressing static non-IID data challenges [CITE:91, CITE:92] to managing dynamic, non-stationary environments. Recent systematic reviews [CITE:2, CITE:33, CITE:79] indicate that traditional FL algorithms degrade when underlying data distributions change over time. To provide a unified perspective, this review categorizes the current literature across three distinct structural dimensions: adaptation strategies (active versus passive), network topologies (centralized, hierarchical, and decentralized), and the primary locus of adaptation (server-side versus client-side).

The first dimension reveals a divergence between active detection and passive adaptation strategies. Active strategies employ explicit drift detection mechanisms—utilizing statistical triggers or performance monitoring—to initiate reactive model updates. For example, frameworks like FL-MalDrift integrate on-device statistical detectors (e.g., ADWIN, DDM) to monitor local data streams and trigger retraining upon detecting distribution shifts [CITE:88, CITE:102]. Similarly, FLARE [CITE:34, CITE:43] utilizes dual-schedulers based on training behavior, while FedConD [CITE:39, CITE:45] and FLAME \cite{Mavromatis2024} incorporate adaptive drift detection and dynamic thresholding. However, these active mechanisms often assume that edge devices can continuously evaluate models on new samples or maintain complex statistical triggers, which can introduce notable computational and latency overheads \cite{ren2025}. 

Conversely, passive approaches eschew explicit detection in favor of continuous adaptation. Algorithms such as Adaptive-FedAVG [CITE:40, CITE:44] utilize dynamic learning rates based on individual training loss to increase model plasticity during drift phases. To address the knowledge erosion (often termed catastrophic forgetting) caused by non-stationary data, Continual Learning (CL) techniques are increasingly integrated into these passive frameworks. This paradigm, Federated Continual Learning (FCL), treats concept drift as a continuous learning problem [CITE:112, CITE:117]. Methods such as CDA-FedAvg [CITE:36, CITE:48] extend standard averaging with continual adaptation techniques, while frameworks incorporating Prioritized Experience Replay or Elastic Weight Consolidation help maintain knowledge across evolving tasks [CITE:118, CITE:128]. By avoiding explicit detection triggers, these passive and FCL mechanisms offer lightweight alternatives that align well with the resource efficiency required in IoT environments.

Beyond adaptation strategies, drift-aware FL architectures vary significantly by network topology. In standard centralized (star) topologies, maintaining a single evolving global model often proves insufficient when clients experience divergent concept drifts. Consequently, multi-model and clustered federated learning approaches, such as FedDrift [CITE:14, CITE:22], FedCCFA [CITE:11, CITE:21], and FedPLC [CITE:20, CITE:37], dynamically group clients and assign specialized models to emerging concepts. While effective for personalized adaptation, these centralized clustering approaches face scalability challenges, including the proliferation of models and complex merging heuristics. To mitigate these bottlenecks, recent literature has explored decentralized (peer-to-peer) and hierarchical topologies. Decentralized FL enables localized drift adaptation through direct client-to-client communication, reducing reliance on a central server and allowing neighborhoods of devices to adapt to regional concept drifts more organically [CITE:142, CITE:155]. Hierarchical models introduce edge servers as intermediaries, aggregating local drifts before global updates, thereby balancing rapid local adaptation with global consistency [CITE:160, CITE:162].

The third dimension involves the locus of adaptation, determining whether the primary responsibility for managing drift resides at the server or the client level. Server-side adaptation relies on the central aggregator to identify specific drifting clients and adjust aggregation weights accordingly. However, because the server is typically restricted from accessing local data distributions due to privacy constraints, it remains partially blind to the exact nature of the drift, complicating its ability to distinguish between drifting and non-drifting clients [CITE:4, CITE:68]. Client-side adaptation, conversely, empowers individual nodes to adjust their local training epochs, learning rates, or regularization terms in response to local data shifts before transmitting updates [CITE:17, CITE:31]. While client-side methods enhance responsiveness, they risk diverging local models so significantly that global aggregation becomes unstable, highlighting an ongoing need for hybrid approaches that balance local plasticity with global stability.

Despite these structural advancements, notable gaps remain in the literature. The intersection of model fairness and distributed concept drift is largely unexplored, presenting an opportunity to develop frameworks that ensure equitable outcomes across sensitive attributes as data distributions evolve [CITE:15, CITE:23, CITE:35]. Furthermore, a primary impediment to advancing drift-aware FL is the absence of unified, standardized benchmarks. Current evaluations frequently rely on fragmented, proprietary datasets, generic data streams, or synthetic adaptations, which limits rigorous cross-paper comparative analyses of non-stationary federated learning algorithms [CITE:54, CITE:81, CITE:122]. Addressing these benchmarking and fairness challenges is essential for the maturation of drift-resilient distributed systems.

\begin{figure}[!t]
\centering
\includegraphics[width=\columnwidth]{figures/top_cited}
\caption{Top cited papers in the survey corpus.}
\label{fig:top_cited}
\end{figure}

\section{Critical Analysis}
\label{sec:critical_analysis}

A critical analysis of the drift-aware federated learning literature reveals several methodological divergences and structural challenges. To establish a unified perspective, this analysis evaluates the state-of-the-art across three foundational dimensions: drift adaptation strategies, the locus of adaptation, and network topologies, followed by an assessment of current benchmarking practices and emerging fairness considerations.

Regarding adaptation strategies, current approaches generally bifurcate into active detection and passive adaptation. Active frameworks, such as FLARE \cite{Chow2023} and FedConD \cite{Chen2021}, utilize explicit drift detection mechanisms—often integrating on-device statistical detectors like ADWIN or DDM to monitor drift scores [CITE:8, 43, 45]. While precise, these active mechanisms frequently assume clients can continuously evaluate models on incoming data streams, which can impose substantial computational and latency overheads in resource-constrained edge and IoT environments [CITE:73, 88]. Conversely, passive adaptation strategies, including Adaptive-FedAVG [CITE:40, 44], bypass explicit detection by dynamically adjusting learning rates based on training loss to enhance model plasticity [CITE:17, 31]. Although more resource-efficient, passive methods must carefully balance plasticity with stability to prevent significant knowledge degradation—often referred to in the literature as catastrophic forgetting. To navigate this stability-plasticity trade-off, Federated Continual Learning (FCL) frameworks increasingly incorporate techniques such as Elastic Weight Consolidation (EWC) and Prioritized Experience Replay (PER) [CITE:48, 112, 117]. However, adapting these methods for distributed environments often introduces additional communication burdens that require careful optimization [CITE:118, 128].

The second dimension, the locus of adaptation, highlights a divergence between client-side and server-side interventions. Client-side adaptation focuses on localized model updates to handle concept shifts independently. In contrast, server-side adaptation typically involves adjusting aggregation protocols or managing model architectures. A prominent server-side debate involves the choice between maintaining a single global model versus deploying multi-model or clustering architectures. Frameworks like FedDrift [CITE:14, 22] and FedCCFA [CITE:11, 21] cluster clients and assign specialized models to distinct emerging concepts [CITE:1, 20]. While effective for divergent distributions, these approaches can encounter scalability limitations over time, potentially leading to elevated storage requirements and complex model-merging procedures. Furthermore, server-side adaptation is inherently constrained by privacy-preserving protocols; the central server lacks direct visibility into local data distributions. Consequently, differentiating between drifting and non-drifting clients during aggregation remains a theoretical challenge, occasionally resulting in suboptimal updates for stable clients [CITE:4, 27].

The third dimension concerns network topologies. The predominant share of the literature focuses on centralized (star-topology) federated networks, where a central server orchestrates drift adaptation. However, there is a notable scarcity of research addressing concept drift within decentralized or peer-to-peer federated learning environments [CITE:55, 62]. In decentralized topologies, the absence of a central aggregator means that local concept drifts propagate asynchronously through peer-to-peer communication. Developing consensus mechanisms that can detect and adapt to non-stationary data streams without a central orchestrator represents a complex, yet underexplored, area of study.

Methodologically, comparative analysis in this domain is constrained by the scarcity of unified, standardized benchmarks for non-stationary federated learning [CITE:94, 122]. Many studies rely on proprietary datasets or simplistic synthetic adaptations, such as sequentially rotating classes in standard computer vision datasets. While useful for initial validation, these synthetic streams do not fully capture the complex, spatio-temporal drift patterns observed in real-world distributed systems, complicating rigorous cross-study evaluations of drift-aware algorithms.

Finally, an emerging area requiring critical attention is the intersection of model fairness and distributed concept drift. As local data distributions evolve, drift can disproportionately impact specific demographic groups or sensitive attributes [CITE:15, 23]. Addressing group-specific distributed concept drift is essential for ensuring that adaptation frameworks maintain equitable performance and do not inadvertently amplify biases in non-stationary environments.

\section{Future Directions and Open Problems}
\label{sec:future_directions}

Building upon the proposed taxonomy—which categorizes the literature by adaptation strategies, network topologies, and the primary locus of adaptation—several underexplored areas remain. Rather than reiterating the mechanistic differences of existing methods, this section synthesizes the survey's core findings to outline actionable research trajectories that will advance the field of distributed non-stationary learning.

**1. Unified Benchmarking Frameworks for Complex Topologies**
A persistent challenge in the current literature is the absence of standardized benchmarks for evaluating federated learning under concept drift \cite{Mahdi2025}. Existing studies frequently rely on fragmented, proprietary datasets or simplistic synthetic adaptations, such as artificially corrupted MNIST variants [CITE:94, 122]. This fragmentation complicates rigorous cross-paper comparative analyses. Future research should prioritize the creation of comprehensive, open-source benchmarking frameworks. Crucially, these benchmarks must accurately simulate real-world distributed non-stationary environments by incorporating temporal drift, spatial non-IID data skew, and support for diverse network topologies (e.g., centralized, hierarchical, and decentralized architectures) [CITE:92, 54, 131].

**2. Topology-Aware Drift Adaptation Strategies**
While the majority of drift-aware federated learning frameworks assume a standard centralized (star) topology, the dynamics of concept drift in alternative network structures remain largely unexplored. In decentralized (peer-to-peer) or hierarchical federated learning, the absence of a global aggregator alters how drift is detected and mitigated [CITE:112, 115]. Future work must investigate topology-aware adaptation strategies, exploring how concept drift propagates through gossip protocols or edge-server hierarchies. Understanding the impact of network topology on drift synchronization and developing localized consensus mechanisms to handle asynchronous concept shifts \cite{Benameur2025} represent critical frontiers for scalable non-stationary learning.

**3. Optimizing the Locus of Adaptation: Client-Side Efficiency vs. Server-Side Attribution**
The literature reveals an ongoing trade-off regarding the primary locus of adaptation. On the client side, active drift detection mechanisms often assume devices can continuously evaluate models on streaming samples, which imposes computational and memory overheads that may be impractical for resource-constrained IoT devices [CITE:34, 88]. While passive adaptation strategies, such as dynamic learning rates, offer lightweight alternatives [CITE:17, 44], they often sacrifice detection precision. Conversely, server-side adaptation faces the challenge of limited visibility; the central server must attribute drift to specific clients without accessing local data distributions. Although some approaches utilize local loss-based detection [CITE:4, 27] or federated fuzzy clustering \cite{Stallmann2024}, these can inadvertently expose client-specific behavioral patterns. Developing hybrid frameworks that balance lightweight client-side heuristics with privacy-preserving, server-side drift attribution [CITE:95, 134] is essential for practical deployment.

**4. Scalable Knowledge Management in Multi-Model Architectures**
To handle distributed concept drift—where different clients experience divergent data distribution shifts—recent works have proposed multi-model and clustered federated learning approaches [CITE:14, 22]. While these methods effectively assign specialized models to handle emerging concepts, they introduce scalability constraints over extended operational periods. Challenges include the unbounded proliferation of concept-specific models, suboptimal model merging, and increased storage and communication overheads [CITE:99, 105]. Future research should focus on scalable multi-model maintenance strategies. Promising directions include leveraging compact model weights for knowledge exchange \cite{Luopan2025}, implementing dynamic repository pruning techniques [CITE:25, 47], and utilizing parameter-efficient fine-tuning (PEFT) to manage concept-specific representations with minimal overhead \cite{paper142}.

**5. Fairness-Aware and Robust Drift Adaptation**
The intersection of model fairness and distributed concept drift presents a vital, yet underexplored, research trajectory. As data distributions evolve, the performance of the global model may degrade disproportionately for certain demographic groups or minority classes, leading to biased outcomes. Addressing "group-specific distributed concept drift" is necessary to maintain equitable performance across sensitive attributes over time [CITE:15, 23]. Future federated adaptation frameworks should incorporate dynamic fairness constraints and adaptive data augmentation techniques \cite{Badar2023} to continuously mitigate discrimination as the underlying data streams evolve, ensuring that drift adaptation does not come at the cost of algorithmic equity.

\section{Conclusion}
\label{sec:conclusion}

This survey has presented a unified perspective on managing concept drift within Federated Learning (FL), moving beyond isolated algorithmic fixes to propose a comprehensive taxonomy. By categorizing the literature across three core dimensions—drift adaptation strategies, network topologies, and the primary locus of adaptation—we have synthesized the current state of drift-aware FL. The analysis indicates a clear transition from localized anomaly detection toward holistic, continuous learning frameworks [CITE:112, 117]. Rather than treating drift as an edge-case disruption, contemporary approaches increasingly embed plasticity directly into the federated optimization process, ensuring models remain relevant across evolving data streams [CITE:122, 128].

Our structured review reveals that effective drift mitigation requires aligning the adaptation strategy with both the network topology and the locus of adaptation. For instance, while centralized FL architectures typically rely on server-side aggregation adjustments to handle global drift [CITE:4, 27], decentralized and hierarchical topologies necessitate client-side or edge-centric mechanisms to manage localized, asynchronous distribution shifts [CITE:14, 22]. Furthermore, the locus of adaptation dictates the distribution of the computational burden. Client-side approaches must balance adaptation efficacy with device constraints [CITE:73, 88], whereas server-side strategies face the challenge of distinguishing between drifting clients and inherent non-IID data heterogeneity without compromising privacy [CITE:39, 45]. 

Despite these advancements, several research avenues require further exploration to mature the field. A primary objective for future work is the development of unified, standardized benchmarks tailored for non-stationary federated environments \cite{Mahdi2025}. The current reliance on disparate datasets and synthetic drift injections limits objective comparative analyses. Establishing open-source frameworks that simulate realistic, spatio-temporal drift across various network topologies will be instrumental in validating emerging algorithms. Additionally, there remains a need for stronger theoretical foundations, particularly in establishing convergence guarantees and generalization bounds for drift-aware aggregation methods under strict privacy constraints [CITE:11, 21].

Scalability also presents an ongoing challenge, particularly as the field explores multi-model and clustered architectures to address divergent client distributions [CITE:15, 23]. Maintaining dynamic model repositories or executing continuous client clustering introduces notable communication and storage overheads [CITE:25, 47]. Future research should investigate more efficient locus-specific adaptations, such as lightweight parameter-efficient fine-tuning (PEFT) at the client level or sparse communication protocols at the server level, to reduce the resource footprint of continuous adaptation mechanisms [CITE:31, 40, 99].

Furthermore, the intersection of model fairness and distributed concept drift warrants deeper investigation. As local data distributions evolve asynchronously, ensuring equitable model performance across diverse demographic groups or client clusters becomes increasingly complex \cite{Badar2023}. Developing fairness-aware adaptation strategies that proactively monitor and mitigate group-specific drift without requiring centralized access to sensitive attributes represents a critical frontier [CITE:15, 44].

In summary, enabling FL to robustly navigate non-stationary environments is essential for its practical deployment in dynamic systems. By aligning adaptation strategies with appropriate topologies and processing loci, and by prioritizing standardized benchmarking, scalable architectures, and fairness, the research community can effectively realize the potential of drift-aware machine learning in distributed settings.

\bibliographystyle{IEEEtran}
\bibliography{references}

\end{document}